% !TEX root = ../main.tex
\lecture{5}{2026-10-19}{Finanzierung}
\textit{siehe Folien} \url{https://lms.uzh.ch/auth/RepositoryEntry/17903616206/CourseNode/114298086656257/path%3D~~Slides~~Unbeschriftet/0} 
\section{Finanzierung}

\subsection{Beteiligungsfinanzierung EK}
\subsubsection{Nennwert vs Aktienwert}

\begin{itemize}
    \item Nennwert = Start des Aktienwertes
    \item \textbf{Berechnung aus Start und Schluss Bilanz}
    \begin{itemize}
        \item Aus 20 Mio Eigenkapital
        \item 200k davon Aktienkapital
        \item $\frac{20 Mio}{200k} = 100 CHF/proaktie$
        \item Nach Gewinne
        \item Eigenkapitalgewachsen 30 Mio/200k Aktienkapital = 150 pro aktie
        \item Nennwert immer noch 100\footnote{noch klären}
    \end{itemize}
\end{itemize}
\subsubsection{Bestandteile des Eigenkapitals}
Art von Aussenfinanzierung des Eigenkapitals
\begin{itemize}
    \item Grundkapital (AK)
    \begin{itemize}
        \item Aktienkapital (inkl. Mitstimmrecht)
        \item Partizipationskapital (nur Beteiligung an Dividende ohne Stimme)
    \end{itemize}
    \item Zuwachskapital (Reserven)
    \begin{itemize}
        \item Offene Reserven
        \begin{itemize}
            \item Kapitalreserven
            \item Gewinnreserven (Pflicht angaben)
            \item Gewinnvortrag (Vorjahres gewinn ausgeschüttet aus beibehalten)
        \end{itemize}
        \item Stille Reserven
    \end{itemize}
\end{itemize}

\subsubsection{Gründe für Kapitalerhöhung}
\begin{itemize}
    \item Finanzierung des Unternehmenswachstums, wenn Kredit und Selbstfinanzierung 
    \item Beseitung einer uNterkapitalisierung (Rekapitalisierung nach verlustjahren)
    \item Ablösung von 
\end{itemize}

\subsubsection{Varianten der Kapitalerhöhung}

\begin{definition}
    [Kapitalerhöhung] das Unternehmen schafft neue Aktien, die anfänglich in ihrem Besitz stehen, und verkaufen diese für einen attraktiveren Preis als die Börsenquotierten Aktien. Die Einahmen aus dem Aktienverkauf sind Finanzierungsquelle
\end{definition}
\begin{enumerate}
    \item Ordentliche Kapitalerhöhung (\textbf{TODO: Vervollständigen})
    \begin{itemize}
        \item Generalversammelung muss beschliessen 
        \item Mit Bezugsrecht
    \end{itemize}
    \item Kapitalband:
    \begin{itemize}
        \item Verwaltungsrat (weniger Menschen als GV) kann kapital für max. 5 Jahre selbständig erhöhen.
    \end{itemize}
    \item Bedingte Kapitalerhöhung
    \begin{itemize}
        \item 
    \end{itemize}
\end{enumerate}

\subsubsection{Bezugsrecht}
\begin{definition}
    [Bezugsrecht] ist Schutzmechanismus für Eigentümer alter Aktien (zu einem höheren Preis gekauft), nach einer Kapitalerhöhung oder Vermehrung der Aktien (mit Minderung des Aktienwerts). Sodass für jede alte Aktien anrecht auf X neue besteht.
\end{definition}

\begin{itemize}
    \item hat einen Wert für den er gekauft werden kann
    \item Aus der Differenz des alten und neuen Aktien preises
    \item schützt vor Verwässerung von:
    \begin{itemize}
        \item Kapital
        \item Gewinn
        \item Stimmrecht
    \end{itemize}
\end{itemize}
\subsection{Kreditfinanzierung FK}

\subsubsection{Formen nach Fristigkeit}

\begin{itemize}
    \item Kurzfristig < 1 Jahr Zahlungsfrist
    \begin{itemize}
        \item Lieferantenkredit
        \item Kundenkredit
        \item Kfr Bankkredit
    \end{itemize}
    \item Mittel fr. (1-3.Jahre) Kapital dass ein Unternehmen braucht um negativen Cashflows zu überbrücken
    \begin{itemize}
        \item Darlehen
        \item Kassaobligation
    \end{itemize}
    \item Lang fr. FK (> 3 Jahre)
    \begin{itemize}
        \item Hypothekardarlehen
        \item Obligationenanleihen
    \end{itemize}
\end{itemize}

\subsubsection{Langfristiges FK: Obligationenanleihen}
\begin{itemize}
    \item Darlehen in Form von Wertpapier in Börse (mit Zahlungsfrist gegen Aktien ohne Frist)
    \item Merkmale
    \begin{itemize}
        \item Betrag (> 50 Mio)
        \item Währung
        \item Stückelung (Aufteilung des gewünschten Betrags in Bonds)
        \item Zins (Meist fix rate pro Jahr)
        \item Laufzeit (3- 10 Jahre)
        \item Ausgabepreis
        \item Rating (Risiko der Anleihe) \footnote{2008 Fälschung von Risiko-Ratings durch Rating-Agenturen}
    \end{itemize}
\end{itemize}
\subsection{Innenfinanzierung}
Wie kommt das Unternehmen selbst dazu zu finanzieren?
\begin{definition}
    [Innenfinanzierung] Generierung von Finanzmitteln aus eigenständiger operativer Tätigkeit (Cashflows)
\end{definition}
\subsubsection{Direkte Ermittlung Cashflows}
Cash-flow = Umsatz - Materialkosten - Personalkosten - Gemeinkosten.

\subsubsection{Indirekte Ermittlung}
Cash-flow  = Reingewinn + Abschreibungen (weil interne Abwertung)

\subsubsection{Abschreibungsgegenwerten}
Die Abschreibung des Maschinenwerts wird am Verkaufpreis der Produkte verteilt und entsprechend erhöht bis es zum Zeitpunkt der Ersatzinvestition kommt.

\textbf{Funktionsweise}
\begin{itemize}
    \item Es entsteht liquidität oder Cashflow obwohl Maschinen nur buchhalterisch Wert verliert

\end{itemize}

\subsubsection{Selbstfinanzierung}
\begin{definition}
    [Selbstfinanzierung] Beschaffung von Kapital durch das Zurückhalten von Gewinnen.
\end{definition}

\textbf{Gründe}
\begin{itemize}
    \item Wichtigste Finanzierungsquelle
    \item Gleichbleibende Beteiligungsverhältnisse ()
    \item Liquiditätsschonend (Keine Zins oder Rückzahlungen nötig)
\end{itemize}
\subsection{Kapitalstruktur}
Fragen: Wie und Warum FK und EK verteilung? Was ist optimale Struktur?

\begin{itemize}
    \item Zieldreick (Rent - Liq - Sec)
    \item Für Rentabilität
    \begin{itemize}
        \item Einsatz von FK kann Redite erhöht werden
        \item Betrag von Reingewinn eigen sinkt aber für investiert EK relativ mehr
    \end{itemize}
\end{itemize}


